\section{Contacts with a path anchored at a random endpoint}
\label{sec:posterior}

The lower contact bound requires a different comparison from the one used in
Section~\ref{sec:entropy}.  A positive bridge now determines the starting
point of the negative path.  We control this dependence by retaining the
conditional probabilities of the bridge's possible endpoints.  The needed
estimate for these probabilities is a finite-vector martingale inequality.

\subsection{The sum of posterior maxima}

The first-moment calculation integrates the nonnegative-martingale maximal
inequality, whose classical form appears in
\cite[Chapter~V, Premi\`ere section, \S2, Theorem~1]{Ville1939}.
The same coordinatewise calculation for posterior vectors appears in
Kesselheim, Molinaro, Patton, and Singla~\cite[Section~1.1]{KesselheimMolinaroPattonSingla2026}.
To obtain an exponential moment, we also control conditional future
increases and the size of each jump. This is the increasing-process
energy mechanism of~\cite[Theorem~4 and its corollary]{Kikuchi1992};
we give an elementary threshold proof with the required filtration explicit.

\begin{lemma}[Posterior maxima]\label{lem:posterior-maxima}
Let $m\ge1$, and let $(w_i(e))_{i\ge0}$, $1\le e\le m$, be nonnegative
martingales for a filtration $(\mathcal F_i)$, with
$\sum_{e=1}^m w_i(e)\le1$ almost surely.  Define
\[
 A_i(e)=\max_{0\le j\le i}w_j(e),\qquad
 A_i=\sum_{e=1}^m A_i(e).
\]
Write $A_\infty(e)=\lim_{i\to\infty}A_i(e)$ and
$A_\infty=\sum_e A_\infty(e)=\lim_{i\to\infty}A_i$.
There are absolute constants $c,C>0$ such that
\begin{equation}\label{eq:posterior-exponential}
 \EE\exp\left(\frac{cA_\infty}{\log(m+1)}\right)\le C.
\end{equation}
Moreover, if $V$ is any event on a joint probability space with this
martingale-path marginal and $p=P(V)$, then
\begin{equation}\label{eq:posterior-event-weight}
 \EE[A_\infty\1_V]
 \le C\log(m+1)\,p\log(e/p),
\end{equation}
where the right side is zero when $p=0$.  The event $V$ need not belong to
any $\mathcal F_i$.
\end{lemma}

\begin{proof}
Write $B=\log(m+1)$.  We first show that, at every almost surely finite
stopping time $\rho$,
\begin{equation}\label{eq:posterior-future-increase}
 \EE[A_\infty-A_\rho\mid\mathcal F_\rho]\le B.
\end{equation}
For a fixed coordinate put $a=A_\rho(e)$ and $p_e=w_\rho(e)$, so that
$0\le p_e\le a\le1$.  Conditional maximal inequality and integration over
levels give, when $a>0$,
\[
 \EE[A_\infty(e)-a\mid\mathcal F_\rho]
 \le\int_a^1\frac{p_e}{u}\,du
 =p_e\log(1/a)\le p_e\log(1/p_e).
\]
If $a=0$, the martingale has conditional expectation zero at every later
time and contributes zero.  The conditional maximal inequality follows
first for bounded stopping times and finite horizons; localization and
monotone convergence give the displayed version.  Summing over $e$ and
adding the missing mass $1-\sum_e p_e$ bounds the resulting sum by the
entropy of a probability vector on $m+1$ elements, hence by $B$.  This
proves~\eqref{eq:posterior-future-increase}.

At each time the increase of a coordinate maximum is at most its current
value.  Consequently
\[
 A_0\le1,\qquad 0\le A_i-A_{i-1}\le\sum_e w_i(e)\le1.
\]
Set $b=2B+1$ and consider the successive thresholds $1+kb$, $k\ge1$.
At their first strict crossings, the overshoot is at most one.  After a
crossing, reaching the next threshold requires a further increase greater
than $b-1=2B$.  Conditional Markov inequality
and~\eqref{eq:posterior-future-increase} therefore bound the conditional
probability of this next crossing by $1/2$.  The first crossing has the
same bound since $A_0\le1$.  Applying this argument at finite horizons
and then letting the horizon increase yields
\[
 P(A_\infty>1+kb)\le2^{-k}.
\]
Because $B\ge\log2$, this is an exponential tail on scale $B$, and its
integral proves~\eqref{eq:posterior-exponential}.

Finally, for $p>0$, conditional Jensen inequality gives
\[
 \exp\left(\frac{c\,\EE[A_\infty\mid V]}{B}\right)
 \le \EE[\exp(cA_\infty/B)\mid V]\le C/p.
\]
Multiplication by $p$ and an adjustment of the absolute constant give
\eqref{eq:posterior-event-weight}.  This last argument uses only a
marginal exponential moment, so it applies to events in a larger joint
space without any assertion that the martingales remain martingales in
an enlarged filtration.
\end{proof}

\subsection{A first-contact estimate}

The comparison of opposing paths through unvisited rows and the use of
quenched exit probabilities are related to Zerner's planar
argument~\cite[Section~2]{Zerner2007}. Here the negative path starts at
a random endpoint of the positive bridge; the posterior estimate controls
the dependence created by that alignment.

We continue to use integer height $x_1$, and write $\pi x\in\ZZ^{d-1}$
for the lateral coordinates of $x\in\ZZ^d$.  Recall that
$u_n^+=P_0(T_n<T_{-1})$ and $p_\pm>0$ are the non-backtracking
probabilities.  The constants below may depend on the environment law and
the dimension.  We assume only strict positivity of the transition rows;
no common ellipticity bound is imposed.

Fix an integer $N\ge3$.  Let $Y$ have the annealed bridge law from $0$ to its first
arrival at height $N$, conditioned on $T_N<T_{-1}$, and put
$E=\pi Y_{T_N}$.  Independently sample a negative path $B$ with law
$\widehat P_{-e_1}$.  Translate it to the random starting site $(N-1,E)$:
\[
 \beta_j=(N-1,E)+B_j,\qquad j\ge0.
\]
Thus the relative negative path is independent of the positive bridge;
its absolute position depends on the bridge endpoint.  A departure of
$Y$ means an index $i<T_N$.  Every site of the infinite path $\beta$ is
a departure site.  For a negative path started at height $N-1$, write
$D^-$ for the event that it never rises above $N-1$.

\begin{proposition}[First contact after a prescribed height]
\label{prop:contact-bound}
In the preceding experiment, let $0<k<r<N$ be integers and define
\[
 \tau=\inf\{i\ge T_k:i<T_N,\ Y_i\in\{\beta_j:j\ge0\}\},
 \qquad
 \Delta=u^+_{r-k}-u^+_{N-k}.
\]
An empty infimum is $+\infty$.  Then
\begin{multline}\label{eq:first-contact-bound}
 P\left(T_r\le\tau<T_N,\
       \min_{T_k\le i\le\tau}(Y_i)_1\ge k,\
       \|E\|\le N^2\right)\\
 \le C\log(2N)\,\Delta\log^2(e/\Delta),
\end{multline}
with right side zero when $\Delta=0$.
\end{proposition}

The estimate is useful when the remaining gap $N-r$ is small relative
to the height $r-k$ already crossed. Lemma~\ref{lem:renewal-tail} gives
$\Delta\le (N-r)F_+(r-k)$; Section~\ref{sec:contradiction} arranges
precisely this separation of scales.

The time $\tau$ refers to the entire negative trajectory and is not
asserted to be a stopping time for the positive path.  We will enumerate
finite prefixes to use its first-contact property.  Martingale estimates
will instead be stopped at deterministic height barriers.

\begin{proof}
Let $\mathcal E_N=\{e\in\ZZ^{d-1}:\|e\|\le N^2\}$.  Its cardinality
satisfies
\begin{equation}\label{eq:endpoint-cardinality}
 \log(|\mathcal E_N|+1)\le C_d\log(2N).
\end{equation}
We first express the event in~\eqref{eq:first-contact-bound} by raw path
weights, then transfer those weights to a shared environment, and only
afterward introduce quenched exit probabilities.

\paragraph{The lower prefix and the endpoint likelihood.}
Enumerate a word $a_0$ from $0$ to its first arrival at height $k$, before
height $-1$.  Its departure sites all have height less than $k$; write
$z$ for its terminal site.  For each $e\in\mathcal E_N$, enumerate a word
$a$ from $z$ to the first contact arrival.  The required words stay in
heights $k,\ldots,N-1$ and reach height $r$ before or at their terminal
time.  Let $\operatorname{Dep}(a)$ denote the departure sites of $a$,
excluding its terminal site unless that site occurred earlier as a
departure.

For any such positive prefix define
\[
 w(a_0a;e)=P_0\bigl(T_N<T_{-1},\ \pi X_{T_N}=e
                   \mid X\text{ begins with }a_0a\bigr).
\]
This is a deterministic function of the two words and $e$.  The two
departure sets of $a_0$ and $a$ are disjoint, so the positive prefix and
completion have raw probability
\begin{equation}\label{eq:prefix-completion-weight}
 W(a_0)W(a)w(a_0a;e).
\end{equation}
For a raw negative walk started at $(N-1,e)$, let $\mathcal B(a)$ be the
event that it never rises above $N-1$, avoids $\operatorname{Dep}(a)$
at every time, and visits the terminal site of $a$.  The first-contact
condition is exactly this avoidance and visit condition.  In particular,
if the terminal site of $a$ occurred earlier as a departure, the event
$\mathcal B(a)$ is empty.  Summing
\eqref{eq:prefix-completion-weight} times the raw probability of
$\mathcal B(a)$, over these words and endpoints, and dividing by
$u_N^+p_-$ gives the probability on the left
of~\eqref{eq:first-contact-bound}.

\paragraph{Transfer and unique prefix counting.}
For a fixed word $a$, the quenched probability of $\mathcal B(a)$ uses
no rows at $\operatorname{Dep}(a)$.  Indeed the walk can be killed on
first arrival at this set, before any row there is used.  This remains
valid for the infinite avoidance event by passage from finite paths.
The departure-row identity of Lemma~\ref{lem:departure-transfer} therefore
gives
\begin{equation}\label{eq:unclassified-lower-transfer}
 W(a)P_{(N-1,e)}(\mathcal B(a))
 =\EE\left[p_\omega(a)
       P_{(N-1,e),\omega}(\mathcal B(a))\right],
\end{equation}
where $p_\omega(a)$ is the quenched product of transition probabilities
along the word.  This identity has not been restricted by any
environment-dependent exit-probability condition.

Keep $W(a_0)$ as a scalar outside the new experiment.  In a fresh
i.i.d. environment sample a raw walk $Z$ from $z$ and a raw negative
walk from $(N-1,e)$, independently conditional on the environment.  We
call $Z$ the comparison path.  Formula~\eqref{eq:unclassified-lower-transfer}
expresses the sum over $a$ as an integral over these two full paths.
For each realized pair, at most one prefix $a$ of $Z$ satisfies the
avoidance and visit conditions: its terminal is the first point of $Z$
in the range of the entire negative path.  This proves uniqueness
pathwise, without using an optional-stopping assertion at contact.
The restrictions that $a$ stays above $k-1$ and reaches $r$ are retained
at this stage.

\paragraph{Stopped posteriors on the comparison path.}
Fix $a_0$ and consider the original raw positive law conditional on this
lower prefix.  In its natural continuation filtration, with no
environment revealed, the coordinates
\[
 w_i(e)=P_0\bigl(T_N<T_{-1},\ \pi X_{T_N}=e
             \mid a_0,\text{ continuation through time }i\bigr),
 \qquad e\in\mathcal E_N,
\]
are nonnegative martingales whose sum is at most one.  Stop them on
first arrival at height $k-1$ or $N$.  All departures up to this stop
have height at least $k$.  These rows are fresh under conditioning on
$a_0$, so the stopped continuation has precisely the same path law as
$Z$ stopped at
\[
 \sigma=T_{k-1}(Z)\wedge T_N(Z).
\]
Arrival at the lower boundary uses a row at height $k$, not a departure
row at $k-1$.  Thus the equality includes the terminal arrival, as
illustrated in Figure~\ref{fig:fresh-strip}.  The stopping time is
almost surely finite by
Lemma~\ref{lem:finite-supremum}: a path that never leaves this finite
height interval would have finite height supremum without tending to
negative height infinity.  The corresponding quenched assertion holds
for almost every environment, simultaneously for all lattice starting
sites and integer height intervals.

Use these same posterior functions on the observed prefixes of $Z$, and
freeze them at $\sigma$.  At arrival at $k-1$ their values need not be
zero: they still describe success before the original lower barrier
$-1$, and that completion may use rows affected by $a_0$.  We neither
recompute them in the fresh environment nor continue them through a
departure below $k$.  Put
\[
 M(e)=\sup_{0\le i\le\sigma}w_i(e),\qquad
 S_*=\sum_{e\in\mathcal E_N}M(e).
\]
Lemma~\ref{lem:posterior-maxima}
and~\eqref{eq:endpoint-cardinality}, applied in the path filtration of
the stopped comparison walk, imply uniformly in $a_0$ that
\begin{equation}\label{eq:comparison-event-weight}
 \EE[S_*\1_V]
 \le C\log(2N)\,P(V)\log(e/P(V))
\end{equation}
for every event $V$ on the joint space of the fresh environment and $Z$.
In particular, $\EE S_*\le C\log(2N)$.  The event-weight inequality,
rather than a martingale assertion after revealing the environment,
permits the next step.

\begin{figure}[tbp]
 \centering
 \begin{tikzpicture}[
  x=1cm,y=1cm,>=Stealth,
  every node/.style={font=\small},
  guide/.style={draw=black!35,densely dashed,line width=.4pt},
  pathline/.style={draw=blue!65!black,line width=1pt}
]
  \fill[black!5] (0,0) rectangle (11.4,1.1);
  \fill[blue!4] (0,1.1) rectangle (11.4,4.6);
  \foreach \yy/\lab in {0/{0},.7/{k-1},1.1/{k},4.6/{N}} {
    \draw[guide] (0,\yy)--(11.4,\yy);
    \node[anchor=east] at (-.15,\yy) {$\lab$};
  }
  \draw[->,black!70] (-.8,.15)--(-.8,4.45);
  \node[rotate=90,anchor=south] at (-1.0,2.3) {height $x_1$};

  \draw[black!65,line width=1pt,->]
    (.65,0)--(.65,.32)--(1.05,.32)--(1.05,.58)
    --(1.45,.58)--(1.45,.88)--(1.75,.88)--(1.75,1.1);
  \fill (1.75,1.1) circle (2pt);
  \node[anchor=north west] at (1.85,1.04) {$z$};
  \node[anchor=west,font=\footnotesize] at (2.65,-.35)
    {lower prefix $a_0$: departures below $k$};

  \draw[pathline,->]
    (1.75,1.1)--(1.75,1.55)--(2.2,1.55)--(2.2,2.1)
    --(2.65,2.1)--(2.65,1.85)--(3.2,1.85)--(3.2,2.55)
    --(3.65,2.55)--(3.65,3.25)--(4.05,3.25)--(4.05,3.9)
    --(4.45,3.9)--(4.45,4.6);
  \draw[pathline,fill=white] (4.45,4.6) circle (2.4pt);
  \node[anchor=south] at (4.45,4.77) {upper exit: stop at arrival};

  \draw[pathline,densely dashed,->]
    (3.2,2.55)--(4.6,2.55)--(4.6,2.15)--(5.0,2.15)
    --(5.0,1.75)--(5.45,1.75)--(5.45,1.1)--(5.9,1.1)
    --(5.9,.7);
  \draw[pathline,fill=white] (5.9,.7) circle (2.4pt);
  \node[anchor=west,align=left,font=\footnotesize,fill=white,inner sep=2pt] at (6.35,.76)
    {lower exit: stop at arrival;\\no departure from height $k-1$};

  \node[anchor=west,align=left] at (6.3,3.6)
    {Rows at heights $k,\ldots,N-1$\\are fresh relative to $a_0$.};
  \node[anchor=west,align=left,font=\footnotesize] at (6.3,2.25)
    {The two exit alternatives are schematic.\\The stopped path laws agree;\\the original posteriors are retained.};
\end{tikzpicture}

 \caption{Fresh rows for the stopped comparison path. The lower prefix
 first reaches height $k$ at $z$, and all its departure rows lie below
 $k$. The continuation, shown with two schematic exit alternatives,
 stops on arrival at height $k-1$ or $N$. It has used only rows with
 heights $k,\ldots,N-1$, so its stopped path law agrees with the fresh
 comparison law. The original endpoint posteriors are frozen at this
 stop, not recomputed under a new completion law. Lateral positions
 are schematic.}
 \label{fig:fresh-strip}
\end{figure}

\paragraph{Quenched exit-probability bins.}
For $k\le y_1<N$ define
\[
 Q_y^\omega=P_{y,\omega}(T_{k-1}<T_N).
\]
Almost surely every transition at every site is strictly positive.
For each interior site $y$, a finite sequence of downward axial steps
reaches $k-1$ before $N$ with positive quenched probability; a finite
sequence of upward axial steps reaches $N$ before $k-1$ with positive
quenched probability.  Thus $0<Q_y^\omega<1$, without a common bound
away from either endpoint.  Partition these values into dyadic bins
$[a,2a)$, where $a=2^{-j}$ and $j\ge1$, and write
$\mathcal D_a^\omega=\{y:k\le y_1<N,\ a\le Q_y^\omega<2a\}$
for the sites in one bin.
Let $V_a$ be the event that, at or after its first arrival at height
$r$ and before $\sigma$, the comparison path visits an interior site with
$Q_y^\omega\ge a$.  With the environment fixed, stop at the first such
visit.  The quenched Markov property gives
\begin{align}
 aP(V_a)
 &\le P_z(T_r<T_{k-1}<T_N)\notag\\
 &=u^+_{r-k}-u^+_{N-k}=\Delta.
 \label{eq:comparison-bin-probability}
\end{align}
The equality follows from nested first-hit events and almost sure exit
from the height interval, again by Lemma~\ref{lem:finite-supremum}.
These events use the fresh comparison law, not the law conditioned on
the lower prefix.

For almost every environment, the raw negative walk has probability at
most $2a$ of both satisfying $D^-$ and hitting $\mathcal D_a^\omega$.
To see this, take its first visit to that set.  On $D^-$ the negative walk
has finite first-coordinate supremum, so Lemma~\ref{lem:finite-supremum} implies that its
first coordinate tends to $-\infty$.  The quenched form of that lemma
holds outside a single environment-null set for all lattice starting
sites.  Consequently, after that first bin visit it must hit $k-1$
before $N$, even if it had previously visited $k-1$.  Its conditional
probability is at most $Q_y^\omega<2a$ by the quenched Markov property.

We now use the unique prefix established above.  In the transferred
integral, classify its contact site by its $Q^\omega$ bin and bound its
endpoint factor $w(a_0a;e)$ by $M(e)$.  Dropping the exact-contact
condition leaves $V_a$ for the comparison path and
$D^-\cap\{\hbox{the negative path hits }\mathcal D_a^\omega\}$
for the negative path.  There is no sum over competing prefixes.
After integrating the negative path, every endpoint $e$ leaves the same
fresh joint marginal $\PP(d\omega)P_{z,\omega}(dZ)$.  Thus the bin's
total contribution, summed over endpoints, is at most
\begin{equation}\label{eq:bin-contribution}
 \begin{split}
 &\sum_{e\in\mathcal E_N}\EE\left[
     M(e)\1_{V_a}
     P_{(N-1,e),\omega}
     (D^-\cap\{\hbox{the path hits }\mathcal D_a^\omega\})\right]\\
 &\hspace{35mm}\le 2a\,\EE\left[\1_{V_a}\sum_{e\in\mathcal E_N}M(e)\right]
  =2a\,\EE[S_*\1_{V_a}].
 \end{split}
\end{equation}

If $\Delta=0$, every $V_a$ is null
by~\eqref{eq:comparison-bin-probability}, proving the required zero
bound.  Suppose $\Delta>0$.  For $a<\Delta$, sum
\eqref{eq:bin-contribution} using $\EE S_*\le C\log(2N)$ and
$\sum_{a<\Delta}a\le2\Delta$.  For $a\ge\Delta$, use
\eqref{eq:comparison-event-weight} and
$P(V_a)\le\Delta/a$ to bound the contribution by
\[
 C\log(2N)\,\Delta\log(ea/\Delta).
\]
There are $O(\log(e/\Delta))$ such dyadic bins, and their logarithms
sum to $O(\log^2(e/\Delta))$.  Thus, for each fixed $a_0$, all endpoint
and upper-prefix contributions together are at most
\[
 C\log(2N)\,\Delta\log^2(e/\Delta).
\]
Finally $\sum_{a_0}W(a_0)=u_k^+\le1$, and
$1/(u_N^+p_-)\le1/(p_+p_-)$.  Restoring these factors proves
\eqref{eq:first-contact-bound}.  The constants lose only these fixed
non-backtracking probabilities and the endpoint-cardinality factor.
In particular, the argument has not required a lower bound for
$Q_y^\omega$ or an inverse-transition moment.
\end{proof}

The estimate controls a contact anywhere after the prescribed height
$r$.  It therefore applies to a long interval without subdividing the
positive path by its successive record heights.  In the next section a
missing interval of contact depths forces exactly such a late first
contact.
